\section*{Chapitre \ref{ch:dv:black-scholes-three-ways} --- Black--Scholes de trois façons}

\begin{solution}{exo:dv:black-scholes-three-ways:1}
Call 7.7603, put 6.2715~; parité call-put~: $7.7603-6.2715=100-100e^{-0.015}$.
\end{solution}

\begin{solution}{exo:dv:black-scholes-three-ways:2}
Pour $V=S$~: $\partial_tV=0$, $\partial_SV=1$, $\partial_{SS}V=0$ et $rS-rS=0$. Pour
$V=e^{rt}$~: $re^{rt}-re^{rt}=0$. Avec $q>0$, $V=Se^{-q(T-t)}$ donne
$qV+(r-q)V-rV=0$~: c'est la valeur d'une action livrée en $T$ sans ses
dividendes, soit le forward prépayé.
\end{solution}

\begin{solution}{exo:dv:black-scholes-three-ways:3}
$0.99\,\bigl(80\,\Phi(d_1)-85\,\Phi(d_2)\bigr)$ avec
$d_{1,2}=\ln(80/85)/(0.35\times0.5)\pm0.0875$~: 3.5572.
\end{solution}

\begin{solution}{exo:dv:black-scholes-three-ways:4}
Le prix du call est le coût de sa réplication avec l'action et le cash~; l'hypothèse de 15~\%
est déjà dans le prix de l'action et modifie les deux membres de la même façon. Le call
ne serait mal évalué que si le trader était en désaccord avec la volatilité qui sera
réalisée~: acheter l'option et pratiquer une couverture en delta rapporte à peu près l'écart entre
la volatilité réalisée et de 20~\%, multipliés par le gamma (chapitre 4). Une opinion sur $\mu$ se
traduit par l'achat de l'action, et non de l'option.
\end{solution}

\begin{solution}{exo:dv:black-scholes-three-ways:5}
$\Phi(d_2)=0.5596$ sous $\mathbb Q$~; avec $\mu=12\%$, $\Phi(0.50)=0.6915$. Le prix
utilise la probabilité risque-neutre, car c'est celle sous laquelle le
coût du \omterm{def:dv:no-arbitrage-and-the-fundamental-theorems:claim}{portefeuille de réplication} est une espérance. Le ratio de couverture vaut $\Phi(d_1)=0.6368$, la
probabilité d'exercice sous la mesure de l'action~: la couverture est la sensibilité au
cours comptant, et non une probabilité de ce qui intéresse le hedger.
\end{solution}

\begin{solution}{exo:dv:black-scholes-three-ways:6}
$-\partial_KC=e^{-rT}\Phi(d_2)$, puisque les termes issus de la dérivation de $d_1$ et $d_2$
se compensent comme pour le delta. Pour les paramètres de référence~: $e^{-0.05}\times0.5596=0.5323$.
\end{solution}

\begin{solution}{exo:dv:black-scholes-three-ways:7}
18.80~\%.
\end{solution}

\begin{solution}{exo:dv:black-scholes-three-ways:8}
L'erreur type Monte-Carlo avec $10^4$ trajectoires vaut 0.148~; 0.19 correspond à 1.3
erreurs types, tout à fait compatible avec la formule. La simulation a utilisé des tirages
lognormaux, elle ne peut donc pas contenir de queues épaisses absentes de la formule~: elle estime le
même nombre, mais avec du bruit. Citer 10.64 reviendrait à citer le bruit.
\end{solution}

\begin{solution}{pb:dv:black-scholes-three-ways:1}
\textbf{1.} Delta 0.6368, gamma 0.0188, theta $-6.414$ par an.
\textbf{2.} $\Theta+rS\Delta+\tfrac12\sigma^2S^2\Gamma-rC=0$ à $10^{-15}$ près.
\textbf{3.} 0.39 (le portefeuille termine à 15.85, l'option paie 15.46).
\textbf{4.} La couverture est quotidienne, non continue~; l'écart est aléatoire, de moyenne proche de
zéro, et il diminue comme l'inverse de la racine carrée du nombre de rééquilibrages.
Le prix est le coût de la stratégie limite.
\textbf{5.} Dans le rendement de l'action, que couvre la vente à découvert~; la
valeur du portefeuille n'en dépend pas.
\textbf{6.} $d_1=0.35$, $d_2=0.15$, $\Phi(d_1)=0.6368$, $\Phi(d_2)=0.5596$.
\textbf{7.} 10.4506 et 5.5735~; $C-P=4.8771=100-100e^{-0.05}$.
\textbf{8.} 0.5596 et 0.6915.
\textbf{9.} $e^{-0.05}\Phi(d_2)=0.5323$.
\textbf{10.} L'approximation donne 8.00. Elle est exacte au premier ordre sur le forward,
105.13, où le call vaut 7.97~; avec un prix d'exercice égal au cours comptant, le call est
dans la monnaie par rapport au forward, de $100-100e^{-0.05}=4.88$ en valeur actuelle.
\textbf{11.} 10.4486, erreur $-0.0020$.
\textbf{12.} 10.4479, erreur type 0.0147, erreur $-0.0026$.
\textbf{13.} Une erreur type de 0.0020 exige $10^6\times(0.0147/0.0020)^2$, soit environ 54
millions de trajectoires.
\textbf{14.} 10.64, erreur type 0.148.
\textbf{15.} La formule pour 10\,000 calls européens (quelques microsecondes chacun)~; un arbre ou une grille
pour un put américain.
\textbf{16.} Volatilité constante~: les volatilités implicites varient selon le strike et l'échéance
(le smile, chapitre 7), et la volatilité réalisée fluctue.
\textbf{17.} La volatilité est l'unité qui permet de comparer des options de strikes, d'échéances et
de sous-jacents différents~; la formule est la conversion, et non une conviction
(livre~1, chapitre~25).
\textbf{18.} Le cours comptant est remplacé par $Se^{-qT}$ dans la formule, et $r-q$ gouverne $d_1$~:
le call tombe à 9.2270.
\textbf{19.} Formule 10.4506~; arbre à 1\,000 pas 10.4486 ($-0.0020$)~; simulation de
Monte-Carlo à $10^6$ trajectoires 10.4479 ($-0.0026$, erreur type 0.0147).
\textbf{20.} Parce que la couverture détient l'action, dont la dérive apparaît alors des deux
côtés de la réplication et se compense.
\end{solution}

\begin{solution}{iq:dv:black-scholes-three-ways:1}
Détenir l'option, vendre $\Delta$ actions à découvert~; la formule d'Itô donne la variation de la position~; choisir
$\Delta=\partial_SV$ pour annuler $dW$~; la position sans risque rapporte $r$~:
$\partial_tV+rS\partial_SV+\tfrac12\sigma^2S^2\partial_{SS}V-rV=0$ (avec $r-q$ dans le
terme de dérive si l'action verse un rendement).
\iqlookfor{l'annulation du terme aléatoire, et l'étape d'absence d'arbitrage.}
\end{solution}

\begin{solution}{iq:dv:black-scholes-three-ways:2}
$0.4\times200\times0.30\times1=24$.
\iqlookfor{la règle $0.4S\sigma\sqrt T$, rapidement.}
\end{solution}

\begin{solution}{iq:dv:black-scholes-three-ways:3}
$\Phi(d_2)$ est la probabilité risque-neutre que le call soit exercé~;
$\Phi(d_1)$ est la probabilité sous la mesure ayant l'action pour numéraire, et
$e^{-qT}\Phi(d_1)$ est le delta. $S\Phi(d_1)e^{-qT}$ est la valeur de la réception de
l'action si l'option est exercée, $Ke^{-rT}\Phi(d_2)$ est la valeur du paiement du prix d'exercice.
\iqlookfor{deux mesures, et la décomposition en digitales asset-or-nothing et
digitales cash-or-nothing.}
\end{solution}

\begin{solution}{iq:dv:black-scholes-three-ways:4}
Parce que l'option est répliquée avec l'action~: le rendement espéré de celle-ci est
déjà inclus dans son prix. Si l'action ne pouvait pas être échangée (une option sur un
indice non négocié, sur une société privée), la réplication échoue et le prix
dépendrait des préférences de risque, donc de la dérive et d'un prix de marché du risque.
\iqlookfor{la réplication comme raison, et l'incomplétude comme exception.}
\end{solution}

\begin{solution}{iq:dv:black-scholes-three-ways:5}
Dynamique lognormale à volatilité constante, négociation continue sans frictions,
taux constants, pas de sauts, un sous-jacent négociable. Les plus coûteuses~: volatilité constante
et l'absence de sauts~; le smile, la volatilité stochastique et les gaps font dépendre le
P\&L de couverture de trajectoires réalisées que le modèle tient pour impossibles.
\iqlookfor{d'abord la volatilité et les sauts, ensuite les coûts.}
\end{solution}

\begin{solution}{iq:dv:black-scholes-three-ways:6}
Vérifier d'abord les bornes~; travailler avec des prix normalisés (diviser par le facteur d'actualisation et
le forward)~; utiliser une fonction d'erreur complémentaire pour les queues~; encadrer et utiliser une
approximation rationnelle comme initialisation, puis un pas de Newton ou de Householder protégé
(la méthode de Jäckel converge en deux itérations)~; vectoriser~; mettre en cache par
échéance~; signaler plutôt que d'écarter les cotations sans solution~; tester les allers-retours sur une grille de
moneyness et maturité à $10^{-12}$ près.
\iqlookfor{bornes, normalisation, queues, garde-fous et grille de tests.}
\end{solution}
